%%=============================================================================
%% CAPITULO 5 -- CONSIDERACOES FINAIS
%%=============================================================================

\section{Considerações finais}

Esta pesquisa busca avaliar como modelos estatísticos espaço-temporais e técnicas de 
aprendizado de máquina aplicados aos dados do SINAPI podem ampliar a capacidade de 
análise e previsão de custos da construção civil. Ao explorar conjuntamente as dimensões 
temporal e territorial dos dados, pretende-se produzir evidências capazes de apoiar o 
planejamento de investimentos públicos e contribuir para decisões relacionadas à alocação 
de recursos em políticas de infraestrutura.

Como contribuição, espera-se demonstrar o potencial de utilização do SINAPI não apenas como 
referência para consulta de preços e elaboração de orçamentos, mas também como fonte de dados 
para análises preditivas aplicadas à gestão pública. Os resultados poderão contribuir para a 
identificação antecipada de tendências de custos, para a avaliação de diferenças territoriais
e para o aprimoramento do planejamento orçamentário, reconhecendo-se que a qualidade das 
previsões dependerá da disponibilidade, consistência e representatividade das séries históricas 
analisadas. 

As limitações e escolhas metodológicas serão consolidadas durante o desenvolvimento 
da pesquisa, especialmente quanto ao recorte temporal, aos algoritmos utilizados e às métricas 
adotadas para comparação dos modelos. Esses elementos também poderão indicar possibilidades 
para pesquisas futuras, incluindo a ampliação da análise para diferentes grupos de insumos, 
localidades e outros fatores capazes de influenciar a evolução dos custos da construção civil.
